\subsection{The two endpoint calculations give a four-volume threshold of four}\label{endpoint-four_volume_threshold--the-two-endpoint-calculations-give-a-four-volume-threshold-of-four}

\subsubsection{Result and motivation}\label{endpoint-four_volume_threshold--result-and-motivation}

The arithmetic endpoint note bounds the monic source norms. The subsequent window-product calculation retains the repeated occurrence of every interior volume contraction. Composing them on the window one degree shorter than the original window improves the original four-volume lower threshold from 2 to 4:

\[
\boxed{\liminf_{k\to\infty}
\frac{\log\bigl(V_{q_k-1}V_{q_k}/(V_{2q_k-1}V_{2q_k})\bigr)}
{q_k\log k}\ge4.}
\]

This statement concerns the original arithmetic family belonging to a fixed packet consisting exactly of a full off-line quartet, with every multiplicity retained. It is a necessary volume-growth consequence for such a packet, not an assertion that an off-line zero exists. No upper bound contradicting this growth is established here.

This note is a written continuation of the two supplied calculations. It is not included in the frozen Toda--Gamma edition DOI \href{https://doi.org/10.5281/zenodo.22730756}{10.5281/zenodo.22730756}.

\subsubsection{1. Original source, constants, and coordinates}\label{endpoint-four_volume_threshold--original-source-constants-and-coordinates}

Keep \(g=2\xi\), the fixed monic packet polynomial \(h\), and

\[
w_h(t)=\frac{|(g/h)(1/2+it)|^2}{2\pi},\quad
m_{h,k}=w_h^{*k},\quad M_h(b)=\int e^{bt}w_h(t)\,dt.
\]

The mass remains \(\mu_h^k\). The monic degree-\(j\) source norm is \(\omega_{h,k,j}\) for the observation \(S=k/2+iu\). In this coordinate map the leading coefficient acquires the phase \(i^j\), of squared modulus one. Neither the source measure nor the generator is replaced.

The arithmetic endpoint proof gives constants \(c_h>0\), \(\vartheta_h=\int_{-1}^1w_h>0\), and \(B_h=42+2\deg h\), such that

\[
m_{h,k}(u)\ge c_h\vartheta_h^{k-3}
 e^{-\pi(|u|+k-3)/2}(1+|u|+k-3)^{-B_h},\qquad k\ge3.
\tag{1}
\]

Here (1) is the proved lower envelope of the actual threefold-and-higher convolution, not a newly postulated comparison measure. Its proof is in the \href{delivery/NOTE.tex}{supplied complete arithmetic note}.

Choose one \(0<b<\pi/2\) and define bounding constants

\[
X_h=\max\{1,M_h(b),M_h(-b)\},\qquad
K_h=\max\{1,\vartheta_h^{-1}\}.
\]

Taking maxima here only bounds constants; it performs no rescaling of the source, its quotient, or either Gram determinant.

\subsubsection{2. The balanced-window norm estimate, with an explicit constant}\label{endpoint-four_volume_threshold--the-balanced-window-norm-estimate-with-an-explicit-constant}

For integers \(n\ge k\ge3\) and \(n/2\le r\le n\), one has

\[
\boxed{\left(\frac{\omega_{h,k,n+r}}{\omega_{h,k,n}}\right)^{1/(2r)}
\le C_h^{\rm bal} n,}
\tag{2}
\]

where the following literal finite constant suffices:

\[
C_h^{\rm bal}=256e^\pi\max\{1,3/c_h\}\,X_hK_h
\max\{1,b^{-3}\}\,6^{B_h/3}.
\tag{3}
\]

To prove this, the trial polynomial \((S-k/2)^j\) and exponential moments give the unchanged upper estimate

\[
\omega_{h,k,j}\le(2j)!b^{-2j}
\bigl(M_h(b)^k+M_h(-b)^k\bigr).
\]

The exact monic Legendre norm on \([-n,n]\) is

\[
\mathfrak l_n(n)=\frac{2n^{2n+1}}{2n+1}
\left(\frac{2^n}{\binom{2n}{n}}\right)^2
\ge\frac{2n^{2n+1}}{(2n+1)4^n}.
\]

This follows from the Legendre leading coefficient and squared norm in \href{https://dlmf.nist.gov/18.3\#T1}{DLMF Table 18.3.1}. Combining (1) with this least-norm calculation gives the exact finite bound

\[
\left(\frac{\omega_{h,k,n+r}}{\omega_{h,k,n}}\right)^{1/(2r)}
\le\left[
\frac{(2(n+r))!b^{-2(n+r)}
\bigl(M_h(b)^k+M_h(-b)^k\bigr)
e^{\pi(n+k-3)/2}(1+n+k-3)^{B_h}}
{c_h\vartheta_h^{k-3}\mathfrak l_n(n)}
\right]^{1/(2r)}.
\tag{4}
\]

For completeness, use \((2(n+r))!\le(4n)^{2(n+r)}\), \(M_h(b)^k+M_h(-b)^k\le2X_h^k\), \(\vartheta_h^{-(k-3)}\le K_h^k\), and \(1+n+k-3\le2n\). After taking the root, the factorial and Legendre factors are bounded by

\[
4n\,8^{n/r}\left(\frac{2n+1}{c_h n}\right)^{1/(2r)}.
\]

The remaining factors are

\[
b^{-(n+r)/r}(X_hK_h)^{k/(2r)}
e^{\pi(n+k-3)/(4r)}(2n)^{B_h/(2r)}.
\]

Now \(n/r\le2\), \((n+r)/r\le3\), \(k/(2r)\le1\), and \((n+k-3)/(4r)\le1\). Also \(2n+1\le3n\), and \(\log(2n)/n\) decreases for \(n\ge3\), so \((2n)^{B_h/(2r)}\le6^{B_h/3}\). These inequalities prove (2)--(3). The diagonal \(r=n\) recovers the scale of the supplied theorem. The needed new case is \(r=n-1\), which lies in the same balanced range.

\subsubsection{3. The exact product on the same canonical metrics}\label{endpoint-four_volume_threshold--the-exact-product-on-the-same-canonical-metrics}

Let \(V_N=\det G_N\), \(\delta_N=V_N/V_{N-1}\in(0,1]\), and retain the calculated real phase \(\phi_N=\sigma-\partial_\theta\log V_N\). The original radius identity is

\[
\epsilon_N^2+\phi_N^2
=\frac{\omega_{N+1}}{\omega_N}
 (1-\delta_N)(\delta_{N+1}^{-1}-1),\qquad N\ge q.
\]

Multiply at \(N=n,\ldots,n+r-1\), without discarding intermediate factors:

\[
\prod_{j=0}^{r-1}(\epsilon_{n+j}^2+\phi_{n+j}^2)
=\frac{\omega_{n+r}V_n}{\omega_nV_{n+r}}
(1-\delta_n)(1-\delta_{n+r})
\prod_{j=1}^{r-1}(1-\delta_{n+j})^2.
\tag{5}
\]

The original phase remains on the left. Each factor on the right after the ratio belongs to \([0,1]\). The exterior theorem bounds the aggregate spectral defect \(L_{h,k}\) by every \(\epsilon_N\). Thus (5) proves

\[
L_{h,k}^{2r}\le
\frac{\omega_{h,k,n+r}V_n}{\omega_{h,k,n}V_{n+r}}.
\tag{6}
\]

This deduction needs neither a new metric nor an assumption that the phase vanishes. A unit contraction makes a factor in (5) zero and directly forces \(L_{h,k}=0\); the inequality (6) still holds. The first admitted degree \(q-1\) is not used in any radius identity here.

\subsubsection{4. Central window and the factor four}\label{endpoint-four_volume_threshold--central-window-and-the-factor-four}

For the exact quartet \(1/2\pm\delta\pm i\gamma\), with \(\delta,\gamma>0\) and common order \(m\), keep

\[
q_k=[1+k(m-1)](k+1)^2,\qquad
L_{h,k}=2\delta[1+k(m-1)](k+1)
\left\lfloor\frac{(k+1)^2}{4}\right\rfloor.
\]

In particular \(L_{h,k}/q_k\ge\delta k/2\). Put \(q=q_k\) and choose \(n=q,r=q-1\). All hypotheses of the proved balanced-window estimate are satisfied: \(q\ge k\ge3\), \(q-1\ge q/2\), and every radius index is at least \(q\). Equations (2) and (6) yield

\[
\boxed{\log\frac{V_q}{V_{2q-1}}
\ge2(q-1)\log\frac{L_{h,k}}{C_h^{\rm bal}q}
\ge2(q-1)\left(\log k+\log\frac{\delta}{2C_h^{\rm bal}}\right).}
\tag{7}
\]

Write \(C_k=\log(V_q/V_{2q-1})\). The original four-volume budget is exactly

\[
\begin{aligned}
\mathcal B_{h,k}
&=\log\frac{V_{q-1}V_q}{V_{2q-1}V_{2q}}\\
&=2C_k+\log\frac{V_{q-1}}{V_q}
              +\log\frac{V_{2q-1}}{V_{2q}}\ge2C_k.
\end{aligned}
\tag{8}
\]

Both last terms are nonnegative by the original rank-one kernel update. Substitution of (7) proves

\[
\boxed{\mathcal B_{h,k}\ge
4(q_k-1)\left(\log k+\log\frac{\delta}{2C_h^{\rm bal}}\right).}
\tag{9}
\]


\subsubsection{The signed metric control at these same four endpoints}
\label{endpoint:metric-backprop}
The AW calculation supplies the full-spectrum estimate and its
condition-number consequence. The SP calculation supplies the additional
actual-overlap refinement. Their exact comparison ACM.1--13 proves that
they act on the same source. We carry both calculations into the budget
(8), rather than replacing its four signed terms by four separate errors.
Here the integration coordinate remains the original $u$ in $S=k/2+iu$;
the coordinate called $y$ in ACM is exactly this $u$. The new variable
$x\in[0,1]$ is only the metric interpolation parameter.

Put $\mathcal P_N=\mathbb C[S]_{\le N}$ and $\mathcal H=\mathcal P_{2q}$ in its original ordered monomials,
$v(u)=(1,S,\ldots,S^{2q})$, and retain the full arithmetic density
$m_1(u)=w_h^{*k}(u)$. The original Gamma comparison has the literal density
and mass
\[
 \begin{gathered}
 r_a(u)=\frac{|\Gamma(a+iu/2)|^2}{2\pi},\qquad
 c_a=2^{1-2a}\Gamma(2a),\\
 m_0(u)=\frac{c_{1/4}^{k}}{c_{k/4}}r_{k/4}(u),\qquad
 c_{1/4}=\sqrt{2\pi},\qquad
 \int_{\mathbb R}m_0(u)\,du=(\sqrt{2\pi})^k.
 \end{gathered}\tag{FV.M1}\label{fvmetric:density}
\]
The complete Gamma convolution proof retained in this reader gives
$m_0=r_{1/4}^{*k}$. Thus (FV.M1) uses its original mass; the arithmetic
mass is still $\mu_h^k$. Define
\[
 m_x=(1-x)m_0+xm_1,\qquad
 M_x=\int_{\mathbb R}v(u)^*v(u)m_x(u)\,du,
 \qquad C_x=M_x^{-1}(M_1-M_0).
 \tag{FV.M2}\label{fvmetric:path}
\]
Both endpoint densities are positive almost everywhere and have all
moments, so every $M_x$ is positive: a nonzero polynomial has only finitely
many zeros on the integration line. No matrix or density is divided by
its mass.

Let $I_N:\mathbb C[S]_{\le N}\hookrightarrow\mathcal H$ be coefficient
inclusion and $B_N:\mathbb C[S]_{\le N-q}\to\mathbb C[S]_{\le N}$ be
multiplication by the original monic $\chi$. Negative-degree spaces are
zero. Put $H_N(x)=I_N^*M_xI_N$, and use the original remainder $J_N$ to set
$G_N(x)=(J_NH_N(x)^{-1}J_N^*)^{-1}$ and $V_N(x)=\det G_N(x)$.
The maps to the common source and their orthogonal projections are exactly
\[
 \begin{aligned}
 \widetilde B_N&=I_N B_N,&
 P_N&=I_N H_N(x)^{-1}I_N^*M_x,\\
 Q_N&=\widetilde B_N
       (\widetilde B_N^*M_x\widetilde B_N)^{-1}
                          \widetilde B_N^*M_x,& Q_{q-1}&=0.
 \end{aligned}\tag{FV.M3}\label{fvmetric:maps}
\]
The formula for $Q_{q-1}$ uses the zero relation domain, with empty
determinant one. For each other $N$, injectivity of multiplication by
monic $\chi$ proves invertibility of its relation Gram. Direct
multiplication proves $P_N^2=P_N$, $Q_N^2=Q_N$ and self-adjointness in
$M_x$, with ranges $\mathcal P_N$ and $\chi\mathcal P_{N-q}$.
These formulas are also the exact typed identities
$\widetilde B_N=B_N^{\rm SP}=I_N B_N^{\rm AW}$.

The coefficient matrix
$[1,S,\ldots,S^{q-1},\chi,S\chi,\ldots,S^{N-q}\chi]$ is triangular
with diagonal one. Its Gram has determinant $\det H_N(x)$.
Eliminating the relation block leaves its Schur complement $G_N(x)$:
indeed for a fixed remainder $z$, minimizing over the relation
coefficients solves the relation normal equations and gives precisely
$z^*G_N(x)z$. Consequently
\[
 \det G_N(x)=\frac{\det H_N(x)}{\det(B_N^*H_N(x)B_N)},\qquad
 \frac{d}{dx}\log V_N(x)=\operatorname{Tr}((P_N-Q_N)C_x).
 \tag{FV.M4}\label{fvmetric:derivative}
\]
For the derivative, multilinearity of determinant gives
$d\log\det H=\operatorname{Tr}(H^{-1}dH)$.
Insert $dH_N=I_N^*(M_1-M_0)I_N$ in the numerator and
$B_N^*dH_NB_N$ in the denominator. Cyclicity of the finite trace
and $M_xC_x=M_1-M_0$ give the two terms in (FV.M4).

Define the original signed function and its two projection sums by
\[
 \begin{aligned}
 F(x)&=\log\frac{V_{q-1}(x)V_q(x)}{V_{2q-1}(x)V_{2q}(x)},\\
 U_x&=Q_{2q-1}+Q_{2q}-Q_q,\\
 W_x&=(P_{2q-1}-P_{q-1})+(P_{2q}-P_q),\qquad A_x=U_x-W_x.
 \end{aligned}\tag{FV.M5}\label{fvmetric:signed}
\]
Keeping the signs of all four terms in (FV.M4) gives
\[
 \begin{aligned}
 F'(x)&=\operatorname{Tr}(A_xC_x)
       =\int_{\mathbb R}(m_1(u)-m_0(u))
                        v(u)A_xM_x^{-1}v(u)^*\,du,\\
 \mathcal B_{h,k}&=F(1)=\mathcal B^\Gamma_{h,k}+\Delta_{h,k},
 \qquad \mathcal B^\Gamma_{h,k}=F(0),\quad
 \Delta_{h,k}=\int_0^1 F'(x)\,dx.
 \end{aligned}\tag{FV.M6}\label{fvmetric:identity}
\]
The scalar density uses the full inverse $M_x^{-1}$, not an ordinary
Euclidean projection kernel. The integral identity follows by substituting
$M_1-M_0=\int v^*v(m_1-m_0)$ in the finite trace; its entries are
integrable moments. The density is real because
$A_xM_x^{-1}$ is Hermitian. It is signed, and its $m_x$ integral
equals $\operatorname{Tr}A_x=0$, as proved next.

The relation subspaces for $q,2q-1,2q$ are nested, of dimensions
$1,q,q+1$. Hence $U_x$ has eigenvalue two on
$\chi\mathcal P_{q-1}\ominus\chi\mathcal P_0$ (dimension $q-1$),
eigenvalue one on $\chi\mathcal P_0$ and
$\chi\mathcal P_q\ominus\chi\mathcal P_{q-1}$ (one dimension each),
and eigenvalue zero on its $q$-dimensional complement.
The degree flag decomposes $W_x$ identically: its two summands have
overlap $\mathcal P_{2q-1}\ominus\mathcal P_q$ of dimension $q-1$,
and two one-dimensional end pieces. Thus both $U_x,W_x$ have rank
$q+1$, trace $2q$ and trace of the square $4q-2$. These statements include
$q=1$, when the overlap dimension $q-1$ is zero.

Let $c_1(x)\le\cdots\le c_{2q+1}(x)$ be the eigenvalues of $C_x$
in $M_x$, let $d_x=c_{2q+1}(x)-c_1(x)$, and put
$s_x=\dim(\operatorname{ran}U_x\cap\operatorname{ran}W_x)$.
Retain both exact pointwise expressions
\[
 \begin{aligned}
 S_{\rm AW}(x)&=c_{q+2}(x)-c_q(x)
       +2\sum_{j=1}^{q-1}(c_{2q+2-j}(x)-c_j(x)),\\
 S_{\rm SP}(x)&=\frac{d_x}{2}\min\left\{4q-2s_x,
 \sqrt{(2q+1)(8q-4-2\operatorname{Tr}(U_xW_x))}\right\},\\
 \mathfrak I_{h,k}&=\int_0^1\min\{S_{\rm AW}(x),S_{\rm SP}(x)\}\,dx.
 \end{aligned}\tag{FV.M7}\label{fvmetric:control}
\]
Here is the full finite proof of their applicability at this earlier
endpoint. For a rank-$r$ orthogonal projection $E$, its diagonal in an
orthonormal eigenbasis of $C_x$ lies in $[0,1]$ and sums to $r$.
Moving this diagonal mass to the lowest or highest eigenvalues gives
$\sum_{j=1}^r c_j\le\operatorname{Tr}(EC_x)
\le\sum_{j=2q+2-r}^{2q+1}c_j$.
Each of $U_x,W_x$ is its range projection, of rank $q+1$, plus its
eigenvalue-two projection, of rank $q-1$. Applying the two inequalities
to these two summands and subtracting the resulting extrema gives
$|F'(x)|\le S_{\rm AW}(x)$, with the multiplicities displayed in
(FV.M7).

Dimension gives $s_x\ge1$. Let $E_x$ be the $M_x$-orthogonal
projection onto the actual intersection of the two ranges. Each of
$U_x,W_x$ dominates its own range projection, which dominates $E_x$;
therefore $U_x-E_x$ and $W_x-E_x$ are positive and each has trace
$2q-s_x$. The triangle inequality for their trace norms gives
$\|A_x\|_1\le4q-2s_x$.
The full spectra just proved also give
\[
 \operatorname{Tr}A_x=0,\qquad
 \operatorname{Tr}A_x^2=8q-4-2\operatorname{Tr}(U_xW_x),\qquad
 \|A_x\|_1^2\le(2q+1)\operatorname{Tr}A_x^2.
 \tag{FV.M8}\label{fvmetric:overlap}
\]
The last inequality is Cauchy--Schwarz on all $2q+1$ real
eigenvalues. Subtract $(c_1+c_{2q+1})I/2$ from $C_x$ in the trace;
the first identity leaves $F'$ unchanged. In an eigenbasis of $A_x$,
each diagonal entry of the centered $C_x$ has absolute value at most
$d_x/2$. Thus $|F'(x)|\le d_x\|A_x\|_1/2\le S_{\rm SP}(x)$.
Both estimates hold at every $x$ on the same original metric. Their
pointwise minimum is integrable: all Grams and inverse Grams are smooth,
all ordered eigenvalues are continuous, and $s_x$ is Borel since its
level sets are determined by vanishing and nonvanishing matrix minors.
Consequently (FV.M6) gives the finite two-sided enclosure
\[
 \boxed{\mathcal B^\Gamma_{h,k}-\mathfrak I_{h,k}
       \le\mathcal B_{h,k}\le
        \mathcal B^\Gamma_{h,k}+\mathfrak I_{h,k}.}
 \tag{FV.M9}\label{fvmetric:budget}
\]

For the positive original relative endomorphism $D=M_0^{-1}M_1$
on $(\mathcal H,M_0)$, let $b_1\le\cdots\le b_{2q+1}$ be all
its eigenvalues. Direct factorization
$M_x=M_0((1-x)I+xD)$ gives
$C_x=((1-x)I+xD)^{-1}(D-I)$ and hence
\[
 \begin{gathered}
 c_j(x)=\frac{b_j-1}{1-x+xb_j},\qquad
 \int_0^1c_j(x)\,dx=\log b_j,\\
 \mathfrak I_{h,k}\le
 \min\left\{\log\frac{b_{q+2}}{b_q}
     +2\sum_{j=1}^{q-1}\log\frac{b_{2q+2-j}}{b_j},\quad
                    \int_0^1S_{\rm SP}(x)\,dx\right\}
 \le(2q-1)\log\frac{b_{2q+1}}{b_1}.
 \end{gathered}\tag{FV.M10}\label{fvmetric:eigenvalues}
\]
The scalar fraction is increasing in $b$ with derivative
$(1-x+xb)^{-2}$, proving the stated eigenvalue order throughout the
path. Integrating its logarithmic derivative proves the integral.
For the last bound use $s_x\ge1$ in (FV.M7) and integrate $d_x$.
The common consequence is the earlier AW bound; the pointwise minimum
retains both its spectral refinement and the additional SP overlap.

The recursive metric transport proof RMT1--22 also applies here on the
actual moment source and supplies a further propagated refinement.
Its coordinate map is $y\mapsto u=y$, leaving $S=k/2+iu$ unchanged.
Equation (FV.M1) equals its reference measure (RMT2)--(RMT3), including
$c_{1/4}^k=(\sqrt{2\pi})^k$; the arithmetic measure is identically
$w_h^{*k}du$ in both. Therefore all entries of its $M_x$ equal (FV.M2).
Its included relation matrix is $I_NB_N=\widetilde B_N$ in (FV.M3),
and its quotient is this same $J_N$. Substitution proves equality of
both Gram inverses, every $P_N,Q_N$, and the four volumes individually.
In particular its $\mathcal B(x)$ equals the original $F(x)$ in (FV.M5),
with no reversal of an endpoint or alteration of a sign. Its quantities
$E_q,O_q,d_x$ are precisely $S_{\rm AW},S_{\rm SP},d_x$ here.

The measure hypothesis of that result holds literally, rather than being
assigned to an arbitrary coefficient form. Its required positivity
$F(x)>0$ can also be checked in these coordinates: if $F(x)=0$,
the two crossed nested ratios $(q-1,2q)$ and $(q,2q-1)$ both have
determinant ratio one. Their least-norm sections, included in the common
$\mathcal H$ by $I_i,I_j$, then agree, because
$G_i-G_j=(I_iR_i-I_jR_j)^*M_x(I_iR_i-I_jR_j)$ is positive and determinant one
forces equality of the two forms. In the first crossed pair the literal
remainder section contains $1$, which would be orthogonal to
$\chi\mathcal P_q$. The polynomial
$\chi^{\#_k}(S)=\sum_j\overline{a_j}(k-S)^j$ for
$\chi(S)=\sum_j a_jS^j$ equals $\overline{\chi(S)}$ on the same
line. Orthogonality to $\chi\chi^{\#_k}$ would force
$\int|\chi(S)|^2m_x(u)du=0$, contrary to the positive degree-$2q$
moment Gram. This retains the full original relation polynomial.

The further two entries in the current recursive control are calculated
on these same operators. Retain
\[
 \begin{aligned}
 T_q(x)&=\frac{d_x}{2}\|U_x-W_x\|_1,\\
 \Sigma_x&=\operatorname{Tr}(C_x^2)
                -\frac{(\operatorname{Tr}C_x)^2}{2q+1},\\
 S_{\rm HS}(x)&=\sqrt{\Sigma_x
                         \operatorname{Tr}((U_x-W_x)^2)}.
 \end{aligned}\tag{FV.M15a}\label{fvmetric:trace-hs}
\]
The midpoint-centering argument proving (FV.M8) already proves
$|F'(x)|\le T_q(x)\le S_{\rm SP}(x)$ before either trace-norm
majorant is taken. Both inequalities are retained. For the second new
entry, $C_x$ is self-adjoint in $M_x$ because $M_xC_x=M_1-M_0$ is
Hermitian. Its centered operator
$C_x-\operatorname{Tr}(C_x)I/(2q+1)$ has squared Hilbert--Schmidt
norm exactly $\Sigma_x\ge0$, by expanding its square and trace.
Since $\operatorname{Tr}(U_x-W_x)=0$, replacing $C_x$ by this
centered operator leaves the trace pairing $F'(x)$ unchanged.
Cauchy--Schwarz on the finite matrix entries in an $M_x$-orthonormal
basis therefore gives $|F'(x)|\le S_{\rm HS}(x)$, with the other
squared norm exactly the trace in (FV.M8). This basis is used only
to prove the Hilbert-space inequality: both endpoint forms and every
source-coordinate operator remain the original ones. The expressions
in (FV.M15a) equal RMT10a--RMT10b after the coefficient identifications
already proved above.

Put $a=2q-1$. The full-angle estimate of RMT11, whose two complete
crossed angle lists and original zero angle are retained in that proof,
therefore gives the simultaneous quantities
\[
 \begin{aligned}
 A_q(x)&=a\sqrt{1-e^{-F(x)/a}}\,d_x,\\
 \mathfrak J_{h,k}&=\int_0^1
       \min\{S_{\rm AW}(x),T_q(x),S_{\rm SP}(x),S_{\rm HS}(x),A_q(x)\}\,dx
                    \le\mathfrak I_{h,k},\\
 \mathfrak I_{h,k}^{\rm nl}&=\int_0^1\min\Bigl\{d_x,
 \frac{S_{\rm AW}(x)}{a\sqrt{1-e^{-F(x)/a}}},
 \frac{T_q(x)}{a\sqrt{1-e^{-F(x)/a}}},\\
 &\hspace{15mm}
 \frac{S_{\rm SP}(x)}{a\sqrt{1-e^{-F(x)/a}}},
 \frac{S_{\rm HS}(x)}{a\sqrt{1-e^{-F(x)/a}}}\Bigr\}\,dx
                    \le\log\frac{b_{2q+1}}{b_1}.
 \end{aligned}\tag{FV.M15}\label{fvmetric:recursive-integrals}
\]
The exact source identities above identify all five retained estimates
individually with RMT9--RMT13, including the exact trace-norm and
centered Hilbert--Schmidt entries. Each bounds the same $|F'(x)|$, so
their pointwise minimum does also. In particular it is at most the
earlier three-entry minimum and the two-entry $\mathfrak I_{h,k}$;
none of their full proofs is discarded. Strict positivity and compactness
of the fixed path keep the displayed denominators nonzero. RMT13 proves
measurability of the actual range intersection, and (FV.M10) computes
the last integral of $d_x$. Integrating gives
$|\Delta_{h,k}|\le\mathfrak J_{h,k}$ directly.

For completeness the scalar coordinate used to propagate this refinement
is $\mathscr H_a(B)=2\operatorname{arcosh}(e^{B/(2a)})$.
It has derivative $1/(a\sqrt{1-e^{-B/a}})$ and inverse
$2a\log\cosh(z/2)$, by direct differentiation and substitution on the
positive real branches. Apply this derivative to the five bounds in
(FV.M15) and integrate, exactly as in the same-source RMT17, to obtain
\[
 |\mathscr H_a(\mathcal B_{h,k})
           -\mathscr H_a(\mathcal B^\Gamma_{h,k})|
       \le\mathfrak I_{h,k}^{\rm nl}.
 \tag{FV.M16}\label{fvmetric:recursive-coordinate}
\]
The constructed signed interval (ISM31)--(ISM37), also propagated in
(RMT19)--(RMT22), is on these identical original forms. Its frequency
$y$ is this $u$, its endomorphism $X_x$ is exactly $C_x=M_x^{-1}(M_1-M_0)$,
and its $\mathcal A_x$ is exactly $U_x-W_x=A_x$ by (FV.M3)--(FV.M6).
Its included relation matrix is $I_NB_N$ and its minimum section is
$I_NR_N(x)$: both have the same remainder, are orthogonal to the same
relation range, and uniqueness of the positive minimum proves equality.
The bare symbol $C:E\to\mathcal H$ in ISM denotes the literal remainder
lift; it is not the endomorphism $C_x$ here. Thus its scalar correction
$\mathcal C_{\rm ISM}=\mathcal B(M_1)-\mathcal B(M_0)$ equals
$\Delta_{h,k}$, with the original Gamma-to-arithmetic orientation.

Fix an integer $p_N\ge0$ solely as a finite truncation order and set
$L_{\rm ISM}=L_{\rm RMT}=p_N$. This symbol changes none of the original
polynomial cutoffs $N$, spectral quantity $L_{h,k}$, or parameters $k,u,x$.
Retain the original positive eigenvalues $b_1,\ldots,b_n$ of
$D=M_0^{-1}M_1$, where $n=2q+1$, and define the exact finite operators
\[
 \begin{aligned}
 B_x^\circ&=(1-x)I+xD,&
 \alpha_x&=\frac{2(1-x)+x(b_1+b_n)}2,\\
 H_x^\circ&=I-B_x^\circ/\alpha_x,&
 C_x^{[p_N]}&=\alpha_x^{-1}\sum_{r=0}^{p_N}(H_x^\circ)^r(D-I),\\
 R_x^{[p_N]}&=(H_x^\circ)^{p_N+1}C_x,&
 E_x^{[p_N]}&=R_x^{[p_N]}-\frac{\operatorname{Tr}R_x^{[p_N]}}n I.
 \end{aligned}\tag{FV.M16a}\label{fvmetric:signed-construction}
\]
Multiplying the finite geometric sum by $I-H_x^\circ$ gives
$C_x-C_x^{[p_N]}=R_x^{[p_N]}$. Every factor is a real function of
$D$ and self-adjoint in the original $M_x$ metric. Trace zero of
$A_x$ therefore gives the exact signed decomposition
$F'(x)=\operatorname{Tr}(C_x^{[p_N]}A_x)+
\operatorname{Tr}(E_x^{[p_N]}A_x)$.
Define the actual center and residual, preserving their signs, by
\[
 \begin{aligned}
 j_{p_N}&=\int_0^1\operatorname{Tr}(C_x^{[p_N]}A_x)\,dx,\\
 e_{p_N}&=\int_0^1
 \sqrt{\operatorname{Tr}((E_x^{[p_N]})^2)\operatorname{Tr}(A_x^2)}\,dx,\\
 \vartheta&=\frac{b_n-b_1}{b_n+b_1},&
 K_D&=\sum_{j=1}^n\left(\frac{|b_j-1|}{\min\{1,b_j\}}\right)^2.
 \end{aligned}\tag{FV.M16b}\label{fvmetric:signed-center}
\]
They are exactly $J_{p_N},\epsilon_{p_N}$ of ISM36 and
$j_{p_N},e_{p_N}$ of RMT20. Cauchy--Schwarz applied to the last
trace in the exact decomposition and then integration proves
\[
 j_{p_N}-e_{p_N}\le\Delta_{h,k}\le j_{p_N}+e_{p_N},\qquad
 0\le e_{p_N}\le\vartheta^{p_N+1}\sqrt{K_D(8q-6)}
                 \longrightarrow0.
 \tag{FV.M16c}\label{fvmetric:signed-error}
\]
For the stated bound, the eigenvalues of $H_x^\circ$ have absolute
value at most $\vartheta<1$, as the midpoint of the extreme eigenvalues
$1-x+xb_1$ and $1-x+xb_n$ is precisely $\alpha_x$.
The original eigenvalues of $C_x$ are $(b_j-1)/(1-x+xb_j)$.
Their absolute values are at most $|b_j-1|/\min\{1,b_j\}$.
Centering the residual subtracts the nonnegative square of its trace
divided by $n$ from its squared Hilbert--Schmidt norm. This proves
$\operatorname{Tr}((E_x^{[p_N]})^2)\le\vartheta^{2p_N+2}K_D$.
Both positive $U_x,W_x$ dominate their range projections, whose
intersection has dimension at least one. Positivity of the trace
product gives $\operatorname{Tr}(U_xW_x)\ge1$, so (FV.M8) gives
$\operatorname{Tr}(A_x^2)\le8q-6$, proving (FV.M16c).
All of these are continuous finite matrix integrands. Their exact
integrals are the objects used here; a numerical integration must
retain its own integration error. Convergence is for this fixed pair
of original source forms. It asserts no estimate uniform in $k$.
The complete original-source and proper-support proof for this
construction is (ISM38)--(ISM43) in
Section~\ref{sec:incoming-source-metric-control}. To compare its primitive
with (FV.M12), instantiate only that primitive calculation with the ordered
pair $(M_1,M_x)$. Its section difference is then exactly
$\rho_{{\rm ISM},1}-\rho_{{\rm ISM},0}
 =I_N(R_N(x)-R_N(1))$, so its boundary has precisely the orientation
in (FV.M12). At $x=0$ this is the negative of the section difference
for the original ordered pair $(M_0,M_1)$; linearity multiplies its
primitive and proper-source residual by $-1$. This typed primitive
comparison leaves the original Gamma-to-arithmetic orientation of
$\Delta_{h,k}$ and the signed estimates above unchanged.

Write $H_0=\mathscr H_a(\mathcal B^\Gamma_{h,k})$ and retain the
complete simultaneous endpoints
\[
 \begin{aligned}
 \mathcal B_{h,k}^{\rm lo}
 &=\max\Bigl\{0,\mathcal B^\Gamma_{h,k}-\mathfrak J_{h,k},
 2a\log\cosh\left(\frac{\max\{0,H_0-\mathfrak I_{h,k}^{\rm nl}\}}2\right),\\
 &\hspace{20mm}\mathcal B^\Gamma_{h,k}+j_{p_N}-e_{p_N}\Bigr\},\\
 \mathcal B_{h,k}^{\rm hi}
 &=\min\Bigl\{\mathcal B^\Gamma_{h,k}+\mathfrak J_{h,k},
 2a\log\cosh\left(\frac{H_0+\mathfrak I_{h,k}^{\rm nl}}2\right),\\
 &\hspace{20mm}\mathcal B^\Gamma_{h,k}+j_{p_N}+e_{p_N}\Bigr\}.
 \end{aligned}\tag{FV.M17}\label{fvmetric:recursive-endpoints}
\]
The inverse is increasing and extends continuously by zero at zero.
Applying it to (FV.M16) and intersecting the resulting interval with
$|\Delta_{h,k}|\le\mathfrak J_{h,k}$ and the oriented interval
(FV.M16c) proves, for every fixed truncation order $p_N\ge0$,
\[
 \mathcal B_{h,k}^{\rm lo}\le\mathcal B_{h,k}
 \le\mathcal B_{h,k}^{\rm hi}
 \le\mathcal B^\Gamma_{h,k}+\mathfrak J_{h,k}
 \le\mathcal B^\Gamma_{h,k}+\mathfrak I_{h,k}.
 \tag{FV.M18}\label{fvmetric:recursive-budget}
\]
These formulas use the complete proof in
Section~\ref{sec:recursive-metric-transport}, with the actual measure
and coefficient correspondences just proved. The preceding generic
signed matrix calculation (FV.M3)--(FV.M10) remains valid on its own
stated domain; the full-angle refinement here is specialized to the
actual original arithmetic/Gamma moment path.

Finally let the two unchanged arithmetic endpoint losses be
$E_0=\log(V_{q-1}/V_q)$ and $E_1=\log(V_{2q-1}/V_{2q})$.
Equation (8), with (FV.M6), gives the literal central identity and bounds
\[
 \begin{gathered}
 2C_k=\mathcal B^\Gamma_{h,k}+\Delta_{h,k}-E_0-E_1,\\
 \max\left\{0,\frac{\mathcal B_{h,k}^{\rm lo}-E_0-E_1}{2}\right\}
 \le C_k\le
 \frac{\mathcal B_{h,k}^{\rm hi}-E_0-E_1}{2},\\
 4(q-1)\log\frac{L_{h,k}}{C_h^{\rm bal}q}+E_0+E_1
 \le\mathcal B_{h,k}\le\mathcal B_{h,k}^{\rm hi}.
 \end{gathered}\tag{FV.M11}\label{fvmetric:central}
\]
The first line is substitution in (8); the second uses (FV.M18) and
$C_k\ge0$; the third uses the first inequality in (7) without
discarding either endpoint loss. All expressions retain the original
full packet $h$, its exact $\chi$, and the constants in (1)--(4).
This gives an actual finite control for the old four-volume expression.
It supplies no uniform estimate of $\mathfrak I_{h,k}$ as the packet
order grows, so the original lower-growth conclusion (9) remains intact.

The comparison is induced by identity polynomial maps on every
$\mathcal P_N$ and the identity on $\mathbb C[S]/(\chi)$.
For the original source realization
$\mathcal V P=P(\sum D_i)F_h^{\otimes k}$ and complete unit map
$\eta[P]=[g/h]_h^{\otimes k}P(\sum s_i)$, these identities commute
exactly with $\eta J_N$; the values and every raw derivative of each
polynomial are unchanged. The least lift at $x$ is
$R_N(x)=H_N(x)^{-1}J_N^*G_N(x)$, so
$J_N(R_N(x)-R_N(1))=0$. Monic division gives a unique coefficient
map $\kappa_N(x)$ with
$R_N(x)-R_N(1)=B_N\kappa_N(x)$.
Retain the original fixed-order tensor division
$\chi(\sum s_i)=\sum_i h(s_i)Q_i(\mathbf s)$. On a quotient vector
$z$, the complete original theta primitive of this difference is
\[
 \sum_{i=1}^k(-1)^{i-1}Q_i(\mathbf D)
       (\kappa_N(x)z)(\textstyle\sum_jD_j)
 \bigl(F_h^{\otimes(i-1)}\otimes\phi_0
                       \otimes F_h^{\otimes(k-i)}\bigr).
 \tag{FV.M12}\label{fvmetric:primitive}
\]
The tensor differential contributes the second sign $(-1)^{i-1}$
and $h(D_i)F_h=\Theta\phi_0$, so its boundary is exactly
$\mathcal V(R_N(x)-R_N(1))z$. Thus the Gram comparison is attached
to an explicit original source map and primitive; it does not identify
the interpolated norm with the original arithmetic norm. For each
unchanged outer label and independent coefficient face, retain the actual
declared subcomplex $C_\lambda^\bullet\subseteq C_{\rm full}^\bullet$
and its degree-$(k-1)$ source $C_\lambda^{k-1}$. Write $\Xi$ for
the primitive in (FV.M12). Whenever $d\Xi$ belongs to $C_\lambda^k$,
its exact residual is $[d\Xi]\in C_\lambda^k/d(C_\lambda^{k-1})$.
The full top-degree comparison has the exact kernel
\[
 \ker\bigl[H^k(C_\lambda^\bullet)\longrightarrow
                 H^k(C_{\rm full}^\bullet)\bigr]
 =\frac{d(C_{\rm full}^{k-1})\cap C_\lambda^k}
             {d(C_\lambda^{k-1})}.
 \tag{FV.M13}\label{fvmetric:proper-source}
\]
Here both complexes have top degree $k$, so all degree-$k$ elements
are cycles. A class $[z]$ maps to zero precisely when $z=d\Xi$ for
some $\Xi\in C_{\rm full}^{k-1}$; it is already zero in the declared
source precisely when $z=d\Xi_\lambda$ for a declared primitive.
This proves (FV.M13), including its injective inclusion into
$H^k(C_\lambda^\bullet)$. Equivalently the primitive map has domain
$\{\Xi\in C_{\rm full}^{k-1}:d\Xi\in C_\lambda^k\}$ and kernel
$C_\lambda^{k-1}+\ker d$, so quotienting by this exact kernel gives
the same residual space. This retains possible cancellations among
tensor slots. For one leg with unchanged top source $\mathscr B$
and injective $\Theta:V\to\mathscr B$, it becomes the exact sequence
\[
 0\longrightarrow V/W\xrightarrow{[v]\mapsto[\Theta v]}
 \mathscr B/\Theta W\longrightarrow\mathscr B/\Theta V
 \longrightarrow0.
 \tag{FV.M14}\label{fvmetric:one-leg-residual}
\]
The first map is injective because $\Theta v\in\Theta W$ forces
$v\in W$; its image is the kernel of the quotient map, which is
surjective by definition. Thus a relation with primitive in the
declared source maps to its supported zero. A general full-source
primitive retains the exact residual (FV.M13) on every coefficient
face, rather than declaring its proper-source class zero. Identity
polynomial maps retain the outer labels; external $\tau$ still maps
to $\tau$ and an empty coefficient face retains its outer label.

Since \(h\) is fixed and \(q_k\to\infty\), division by \(q_k\log k\) gives the claimed liminf of at least 4. An upper limit strictly below 4 would already conflict with (9); no such upper estimate has been obtained here. The two-volume conclusion at \(n=r=q\) is also retained: \(\liminf \log(V_q/V_{2q})/(q\log k)\ge2\).

\subsubsection{5. Source maps and verification scope}\label{endpoint-four_volume_threshold--source-maps-and-verification-scope}

All representatives remain those of the supplied arithmetic construction: \(J^{(k)}R_N=\eta_{h,k}\) and \(q^{(k)}R_N=\sigma_h^{\otimes k}\eta_{h,k}\). Restricting the degree window does not alter their unit, jets, fibre labels, or cohomology classes. The supported-zero relation and the external absence \(\tau\) retain their original separate maps. In particular, no repeated exterior iteration or erased multiplicity cost occurs in (5)--(9).

The analytic envelope and original norm theorem are from the supplied Arithmetic Endpoint Bounds note. Equation (5) is the subsequent supplied window-product calculation. Equations (2)--(4) and their composition (7)--(9) are the present continuation. The two independent review lanes record finite tests separately from written analytic proof; no new Lean certificate or interval enclosure of \(C_h^{\rm bal}\) is asserted.
